\subsection{Experimental setup}
\paragraph{Real LLM and benchmark.}
\emph{Data.} MMLU-Pro \citep{wang2024mmlupro} (dataset revision \texttt{b189ec7}, 12\,032 test questions, up to 10 options). We fix a uniformly random population $P$ of $\RealN$ test questions; the remaining questions are never used. The 70 validation questions provide few-shot demonstrations.
\emph{Model.} Qwen2.5-1.5B-Instruct \citep{qwen2024qwen25}, 8-bit GGUF, run with llama.cpp on CPU; the answer is the option letter with the largest next-token log-probability (greedy, deterministic, no chain of thought), so a (configuration, question) pair has one fixed outcome.
\emph{Configuration space} ($K{=}72$): option format (\texttt{A.\ text} / \texttt{(A) text}) $\times$ demonstrations (none / one of two fixed 3-shot sets) $\times$ one of 4 instruction sentences $\times$ one of 3 answer cues. This is a small, hand-built space of the kind that instruction/demonstration optimisers explore, evaluated exhaustively so that every search is a replay on a complete per-item score matrix $Y\in\{0,1\}^{|P|\times 72}$.
\emph{Benchmark draws and truth.} Each replicate draws a benchmark $S$ of $n\in\{50,100,200,400,800\}$ questions i.i.d.\ (with replacement) from $P$, and the true performance of a configuration is its exact accuracy over $P$: the empirical distribution of $P$ plays the role of the task distribution. (Drawing $S$ without replacement and scoring on $P\setminus S$ is also possible, but then the fresh accuracy is mechanically anti-correlated with the benchmark accuracy and the measured optimism is inflated by $|P|/(|P|-n)$; App.~\ref{app:wr} reports this secondary design for small $n$.) \emph{Search procedures} (each sees only $S$): random search with budget $B\in\{1,\dots,72\}$ and full evaluation of each candidate; successive halving (3-way, nested subsets of $S$, final rung on all of $S$); greedy coordinate ascent over the three factors from a random start; and a surrogate-guided search in the style of MIPRO \citep{opsahl2024optimizing}: 40 trials on random minibatches of 50 questions proposed from an additive per-factor surrogate, then full evaluation of the 3 best. These are our own implementations of the algorithm families, \emph{not} the DSPy code (Sec.~\ref{sec:disc}).

\paragraph{Synthetic worlds.}
To control the quantities that theory says matter and to reach regimes the real pool does not contain, we simulate item-level outcomes from a latent-threshold model with the same factorial configuration structure: $y_{ci}=\mathbb 1[a+\mu_c+\sqrt{\rho}\,z_i+\sqrt{1-\rho}\,u_{ci}>0]$, $z_i,u_{ci}\sim\mathcal N(0,1)$, where $\mu_c$ (standardised to SD $\tau$) mixes additive factor effects and interactions, $\rho$ is the latent between-configuration correlation (item difficulty shared by all configurations), and $a$ sets base accuracy. We also use i.i.d.\ configuration effects with normal, sparse (5\% of configurations much better), heavy-tailed ($t_3$) and bimodal (half the configurations ``broken'') distributions for robustness checks.

\paragraph{Metrics.}
For each reporting method: \emph{bias} $=\mathbb{E}[\text{report}-\mu(\hat\theta)]$, \emph{RMSE}, \emph{coverage} and median width of nominal-95\% intervals for $\mu(\hat\theta)$ (the quality of the configuration that would actually be deployed), and, for comparisons, the rate of significant ``$A>B$'' claims at one-sided level 5\% against the true gap. Monte-Carlo errors are reported (bias: $\pm1.96$ SE; coverage: binomial SE).
